\section{Introduction}
\label{sec:intro}

A semantic cache answers a new prompt with a stored response whenever the new prompt
looks close enough to one the system has already served. Closeness is judged in an
embedding space: the cache encodes the incoming query, finds its nearest stored
neighbor, and returns the neighbor's answer when their cosine similarity clears a
threshold $\thr$. The appeal is obvious. A served hit skips a model call, so a cache
that fires often turns the dominant cost of an \mbox{LLM} service, inference, into a
lookup~\cite{bang2023gptcache,gill2024meancache}. Production systems now ship semantic
caches by default, and recent work has pushed their hit rates higher and added online
guarantees on correctness~\cite{schroeder2025vcache,wang2025categoryaware}.

The threshold hides a hard decision. Set it low and the cache fires often, paying back
its keep, but it also begins to serve answers to prompts that only resemble an earlier
one. Set it high and those mistakes vanish, along with most of the savings. The mistake
itself is quiet. When the cache returns a stored answer to a prompt it does not actually
fit, nothing errors and nothing logs; the user simply receives a confident, wrong
reply. We call this a \emph{false hit}, and it is the failure mode this paper measures.

The difficulty is that embedding similarity and answer reuse are not the same property.
Two prompts can be near-identical in wording yet require opposite answers, as ``what are
the strengths of this design'' and ``what are the weaknesses of this design'' do; two
others can share almost no words yet take the same answer. Similarity is what the cache
can see. \emph{Cache-equivalence}, the question of whether one prompt's stored answer is
valid for another, is what it needs. Prior work has largely optimized the first as a
proxy for the second: it tunes thresholds, learns them online from correctness
feedback~\cite{schroeder2025vcache}, specializes them by query
category~\cite{wang2025categoryaware}, or sharpens the embedding
itself~\cite{gill2025advancing}. What the literature has not done is measure how often
the proxy fails when it is not under attack, chart how that failure rate moves with the
embedding model and the domain, follow it into the quality of the task the cache serves,
or turn the measurement into an operating point that accounts for what a false hit
actually costs. The one adversarial study of these caches shows the proxy can be broken
on purpose~\cite{zhang2026keycollision}; we ask how often it breaks on its own.

This paper answers that question with a measurement study and a calibration method built
on it. We separate cache-equivalence from textual similarity and label it directly,
drawing on human-annotated pair corpora whose judgments encode same-meaning rather than
word overlap~\cite{zhang2019paws,wang2019glue,dolan2005mrpc}. On this footing we measure
the false-hit rate of six widely used sentence
embeddings~\cite{reimers2019sentencebert,wang2022e5,xiao2023cpack} across three domains
and the full range of thresholds, report it with bootstrap confidence intervals, and
read off the operating points a practitioner reasons about: how much of the traffic a
cache can serve while holding its error below a ceiling, and how clean it can be while
still reusing a target share of genuinely reusable answers. We then inject false hits
into downstream retrieval and generation to measure what they cost once served, and we
use that cost to calibrate the threshold. The resulting policy minimizes expected cost
under the measured, asymmetric penalty of a false hit, and it transfers across models
and domains without the per-prompt online feedback that earlier correctness-bounding
methods require.

Our findings are not reassuring for the default of trusting a single similarity threshold.
On adversarial paraphrases the strongest embedding reaches a \rocauc{} of only $0.67$, and a
lexical baseline outscores every neural encoder, because the surface-form insensitivity that
makes an embedding useful is exactly what blinds it to a word-scrambled near-duplicate.
Holding the false-hit rate under one percent caps coverage near nine percent on the easiest
domain and under one percent on the others. Downstream, a false hit is nearly as costly when it
corrupts the retrieved context as when it corrupts the final answer, because the reader rarely
recovers from a wrong premise; what softens the damage is a milder corruption, not an earlier
position in the pipeline. A fixed similarity threshold proves
catastrophic under a realistic false-hit penalty, costing several times the price of not
caching at all, while a threshold calibrated to the measured penalty recovers near-baseline
cost and adapts to each domain's difficulty. Such a threshold transfers only in the
conservative direction: a value tuned on an easy domain backfires when reused on a hard one.

\paragraph{Contributions.}
\begin{itemize}
  \item A \emph{cache-equivalence} construct that separates ``safe to reuse the stored
        answer'' from ``looks similar,'' together with a labeling protocol that grounds
        it in human judgments rather than embedding geometry (\cref{sec:setup}).
  \item The first systematic map of benign false-hit economics for embedding caches:
        false-hit rate and coverage across six encoders, three domains, and all
        thresholds, with bootstrap confidence intervals and the operating points that
        matter in deployment (\cref{sec:results}).
  \item A measurement of the downstream cost of a false hit, obtained by controlled
        fault injection into retrieval and generation, broken down by task and by the
        cache's position in the pipeline (\cref{sec:downstream}).
  \item A cost-aware threshold-calibration policy that ties the operating point to the
        measured cost of a false hit. It exposes that a fixed default is catastrophic under a
        realistic penalty, matches a correctness-bounding baseline without a hand-set error
        target, and reveals that thresholds transfer across domains only in the conservative
        direction (\cref{sec:calibration}).
\end{itemize}

All code, the labeled data, and the exact commands that regenerate every number and
figure are released with the paper.
